\section{Related Work}
\label{sec:related}

\paragraph{Universal adversarial attacks on multimodal LLMs.}
The Universal Adversarial Attack (UAA) of \citet{rahmatullaev2025universal} optimises a single image that, when paired with any of $60$ benign prompts, drives an aligned multimodal LLM toward a chosen target phrase. They report jailbreak ASR up to $81\%$ on SafeBench / MM-SafetyBench. Earlier, \citet{schlarmann2023robustness} showed that perturbations as small as $\eps=1/255$ can re-route VLM captions to attacker-chosen URLs, and \citet{qi2024visual} demonstrated that one visual adversarial example can universally jailbreak aligned VLMs into following harmful instructions outside the optimisation corpus. \citet{carlini2024aligned} further established that multimodal models are roughly an order of magnitude easier to break than their text-only counterparts. \citet{bailey2024image} introduced ``image hijacks'' --- behaviour-matching adversarial images for output-control, context-exfiltration, safety-override, and false-belief attacks --- with ASR $\geq 80\%$ against LLaVA. We build on \citet{rahmatullaev2025universal} directly: their PGD optimisation is our Stage~1 (\S\ref{sec:method-stage1}). Our contribution relative to these works is the post-fusion dual-axis evaluation, which separates the broad-but-shallow disruption their methods deliver from the rare-but-targeted payload delivery they sometimes claim.

\paragraph{AnyAttack and CLIP-based fusion.}
\citet{zhang2025anyattack} train a self-supervised CLIP-encoder + decoder pair on bidirectional COCO that maps any image's CLIP feature to an \Linf{}-bounded noise tensor. The decoder is general-purpose: it produces transferable adversarial perturbations against frontier closed VLMs without per-target retraining. We reuse the public \texttt{coco\_bi.pt} weights as our Stage~2 (\S\ref{sec:method-stage2}). \citet{zhang2024adversarial}'s Adversarial Illusions in Multi-Modal Embeddings (USENIX Security 2024 distinguished paper) is a closely-related embedding-space attack against ImageBind / AudioCLIP. \textsc{AttackVLM} \citep{zhao2023attackvlm} is the de-facto baseline for ``targeted black-box attack on open VLMs'' --- transferred from CLIP / BLIP to MiniGPT-4, LLaVA, BLIP-2.

\paragraph{Indirect prompt injection.}
\citet{greshake2023indirect} (AISec'23) introduced the indirect-prompt-injection (IPI) taxonomy --- attacks that smuggle instructions through retrieved data or third-party documents rather than directly through the user prompt. \citet{bagdasaryan2023imagessounds} ported this idea to images and audio, demonstrating adversarial perturbations that act as instruction injections against LLaVA and PandaGPT. \citet{liu2024formalizing} (USENIX Security 2024) provide a formal definition --- a successful injection requires the target string to appear in the response --- which we adopt as the basis of our Target Injected check. \citet{yi2025bipia} (KDD 2025) build the first IPI benchmark for text LLMs (BIPIA) with boundary-awareness defenses; \citet{zhan2024injecagent} (ACL Findings 2024) extend it to tool-using agents. Our work is the visual-modality analogue of these benchmarks, with explicit dual-axis scoring.

\paragraph{Multimodal jailbreaks against aligned VLMs.}
\citet{zou2023gcg} introduced GCG --- a transferable suffix attack on aligned LLMs that became the textual analogue of universal visual attacks. The visual branch has produced a growing list of methods: \textsc{Jailbreak-in-Pieces} \citep{shayegani2024jailbreak} (ICLR 2024 Spotlight) splits the payload across image and text modalities; \textsc{FigStep} \citep{gong2025figstep} (AAAI 2025 Oral) renders prohibited text typographically as $82.5\%$ ASR average across six open VLMs; \textsc{HADES} \citep{li2024hades} (ECCV 2024 Oral) combines typography with diffusion-generated harmful imagery and adversarial perturbation, hitting $90.3\%$ ASR on LLaVA-1.5 and $71.6\%$ on Gemini Pro Vision; \textsc{BAP} \citep{ying2024bap} jointly optimises image and text prompts for a $+29\%$ boost over single-modality. \textsc{ImgTrojan} \citep{tao2025imgtrojan} (NAACL 2025) shows that data-poisoning at $0.0001\%$ ratio (1/9{,}198) suffices to jailbreak LLaVA-v1.5 (orthogonal threat model: training-time vs. inference-time). We differ from this entire family in threat model: we attack \emph{output integrity} under \emph{benign user prompts}, whereas these works attack safety alignment under adversarial prompts. The high-ASR numbers in this literature partly motivate our methodological critique --- they measure ``safety failure,'' not ``payload delivery,'' and the two are not the same.

\paragraph{Adversarial robustness and defenses.}
\citet{schlarmann2024robustclip} (ICML 2024 Oral) propose Robust CLIP --- an unsupervised adversarial fine-tuning of the vision encoder that drops in to any LVLM. \citet{hossain2024simclip} extend this with Sim-CLIP. \citet{zong2024vlguard} (ICML 2024) introduce VLGuard, a safety-tuning dataset that brings adversarial ASR near zero in many cases. \textsc{ECSO} \citep{gou2024ecso} (ECCV 2024) routes suspect images through OCR/captioning, falling back to the text-only safety pipeline --- a $+37.6\%$ MM-SafetyBench improvement, and the defense most likely to neutralise our pipeline. \textsc{AdaShield} \citep{wang2024adashield} (ECCV 2024) provides adaptive defense prompts. \textsc{DiffPure} \citep{nie2022diffpure} (ICML 2022) removes adversarial noise via diffusion forward+reverse. We do not propose a defense in this paper; \S\ref{sec:discussion} discusses which defenses our findings predict to be effective. Foundational adversarial-attack methodology comes from \citet{goodfellow2015explaining} (FGSM), \citet{madry2018pgd} (PGD adversarial training), and \citet{moosavi2017uap} (universal adversarial perturbations).

\paragraph{Evaluation methodology and benchmarks.}
\textsc{HarmBench} \citep{mazeika2024harmbench} (ICML 2024) standardises evaluation across $18$ attacks $\times$ $33$ LLM/defenses with four behaviour categories. \textsc{JailbreakBench} \citep{chao2024jailbreakbench} (NeurIPS 2024 D\&B) provides an open evolving repo of $100$ behaviours with leaderboard scoring. \textsc{MM-SafetyBench} \citep{liu2024mmsafetybench} (ECCV 2024) extends to the multimodal setting with $5{,}040$ image-text pairs across $13$ unsafe scenarios. \textsc{MMJ-Bench} \citep{weng2025mmjbench} (AAAI 2025) is the first unified VLM-jailbreak attack-defense leaderboard. \textsc{RTVLM} \citep{li2024rtvlm} covers $10$ VLM safety subtasks. None of these benchmarks separate Output Affected from Target Injected; their scoring is a single Boolean ``did the model produce harmful / target content'' typically judged by an LLM-as-judge \citep{zheng2023judge}. Our contribution is to expose this gap by reporting both numbers explicitly across $6{,}615$ pairs --- we find a $290\times$ divergence on the same data. The released dataset (\S\ref{sec:conclusion}) is intended to slot into HarmBench / JailbreakBench-style regression suites with both axes available per pair.
